% Part: second-order-logic
% Chapter: syntax-and-semantics
% Section: semantic-notions

\documentclass[../../../include/open-logic-section]{subfiles}

\begin{document}

\olfileid{sol}{syn}{sem}
\olsection{معنوی تصورات}

\begin{explain}
\emph{منطقی اعتبار}، \emph{منطقی نتیجہ} تے
\emph{پورا ہون جوگتا} دیاں تعریفاں دوجے درجے دی منطق وچ
اوہو جیہیاں نیں جیہڑیاں پہلے درجے دی منطق وچ نیں۔ فرق ایہ
اے کہ ہُن پورا اترن دا تعلق دوجے درجے دے !!{formula}s لئی
ورتیدا اے۔ دوجے درجے دا
!!{sentence} اوہ !!a{formula} اے جیہدے سارے
متغیر پابند ہون، محمولی تے فنکشنی متغیر وی۔
\end{explain}

\begin{defn}[منطقی اعتبار]
جملہ $!A$ \emph{معتبر} اے، $\Entails !A$،
اُدوں تے صرف اُدوں جدوں $\Sat{M}{!A}$ ہر
!!{structure}~$\Struct M$ لئی ہووے۔
\end{defn}

\begin{defn}[منطقی نتیجہ]
جملیاں دے سیٹ~$\Gamma$ توں جملہ~$!A$
\emph{منطقی طور تے نکلدا اے}، $\Gamma
\Entails !A$، اُدوں تے صرف اُدوں جدوں ہر اوہ !!{structure}~$\Struct M$ جیہدے وچ
$\Sat{M}{\Gamma}$ ہووے، اوہدے وچ $\Sat{M}{!A}$ وی ہووے۔
\end{defn}

\begin{defn}[پورا ہون جوگتا]
جملیاں دا سیٹ~$\Gamma$ \emph{پورا ہون جوگا} اے
جے $\Sat{M}{\Gamma}$ کسے
!!{structure}~$\Struct M$ لئی ہووے۔ جے
$\Gamma$ پورا ہون جوگا نہ ہووے، تاں اوہنوں
\emph{پورا ہون جوگا نہیں} آکھدے نیں۔
\end{defn}

\end{document}
